\documentclass[12pt,a4paper]{article}

\usepackage[utf8]{inputenc}
\usepackage[T1]{fontenc}
\usepackage[brazil]{babel}
\usepackage{geometry}
\usepackage{graphicx}
\usepackage{hyperref}
\usepackage{longtable}
\usepackage{array}
\usepackage{xcolor}
\usepackage{enumitem}

\geometry{margin=2.5cm}
\hypersetup{
    colorlinks=true,
    linkcolor=blue,
    urlcolor=blue,
    citecolor=blue
}

\title{Manual do Sistema Quintana}
\author{Equipe do Projeto}
\date{\today}

\begin{document}

\maketitle

\begin{center}
\textbf{Versão do documento:} preencher \\
\textbf{Versão/commit do sistema:} preencher \\
\textbf{Responsáveis pela documentação:} preencher
\end{center}

\newpage
\tableofcontents
\newpage

\section{Apresentação}

\subsection{Objetivo do sistema}
Descrever, em linguagem acessível, o que é o Quintana e qual problema educacional ele apoia.

\subsection{Público-alvo}
Indicar os perfis de usuário atendidos:
\begin{itemize}
    \item estudantes;
    \item professores;
    \item equipe responsável pela oficina ou operação.
\end{itemize}

\subsection{Escopo do manual}
Explicar o que este manual cobre e o que não cobre.

\subsection{Convenções usadas no documento}
Explicar padrões de escrita, nomes de botões, exemplos de usuários, comandos e imagens.

\section{Visão geral do sistema}

\subsection{Principais funcionalidades}
Listar as funcionalidades centrais do sistema.

\subsection{Fluxo geral de uso}
Inserir um diagrama ou descrição textual do fluxo:
\begin{enumerate}
    \item login;
    \item escolha do perfil;
    \item envio ou consulta de redações;
    \item visualização de resultados;
    \item análise pedagógica.
\end{enumerate}

\subsection{Perfis de usuário}

\subsubsection{Estudante}
Descrever o que o estudante pode fazer.

\subsubsection{Professor}
Descrever o que o professor pode fazer.

\section{Perfis e autenticação}

\subsection{Endereços e ambientes}
Informar as URLs usadas em ambiente local, oficina ou produção.

\subsection{Modelo de autenticação}
Descrever, em nível técnico-operacional, como a autenticação funciona:
\begin{itemize}
    \item login por e-mail e senha;
    \item emissão de token;
    \item envio do token nas chamadas autenticadas;
    \item expiração do token.
\end{itemize}

\subsection{Usuários fictícios para demonstração}
Preencher tabela com usuários de teste.

\begin{longtable}{|p{4cm}|p{4cm}|p{4cm}|}
\hline
\textbf{Perfil} & \textbf{E-mail} & \textbf{Senha} \\
\hline
Professor & preencher & preencher \\
\hline
Aluno & preencher & preencher \\
\hline
\end{longtable}

\subsection{Perfis e permissões}
Resumir as permissões de estudante e professor.

\section{Funcionalidades por perfil}

\subsection{Observação sobre o manual de utilização}
O passo a passo de uso das telas será descrito em um manual de utilização separado.
Neste documento, as funcionalidades devem ser descritas em nível de sistema: objetivo, dados envolvidos, regras principais e dependências.

\subsection{Funcionalidades do estudante}
Descrever, de forma resumida, os recursos disponíveis ao estudante:
\begin{itemize}
    \item consulta de temas;
    \item envio de redações;
    \item visualização de notas por competência;
    \item radar médio das competências;
    \item linha do tempo de progresso;
    \item plano de ação;
    \item checklist de reescrita;
    \item comparação entre versões;
    \item atividades atribuídas.
\end{itemize}

\subsection{Funcionalidades do professor}
Descrever, de forma resumida, os recursos disponíveis ao professor:
\begin{itemize}
    \item cadastro e gestão de temas;
    \item consulta de redações recebidas;
    \item correção ou comentário do professor;
    \item cadastro e gestão de turmas;
    \item cadastro e gestão de atividades;
    \item acompanhamento de submissões;
    \item análise da turma;
    \item alertas pedagógicos.
\end{itemize}

\subsection{Relação com o manual de utilização}
Indicar quais tópicos deverão ser detalhados no manual de utilização, evitando duplicação neste documento.

\section{Arquitetura funcional}

\subsection{Módulos do frontend}
Descrever as principais páginas e componentes do frontend, sem detalhar o passo a passo de uso.

\subsection{Módulos do backend}
Descrever os principais módulos do backend:
\begin{itemize}
    \item autenticação e autorização;
    \item avaliação de redações;
    \item feedback estruturado;
    \item analytics do professor;
    \item acesso ao banco;
    \item scripts operacionais.
\end{itemize}

\subsection{Coleções do MongoDB}
Descrever as coleções principais e seus relacionamentos:
\begin{itemize}
    \item users;
    \item temas;
    \item classes;
    \item activities;
    \item redacoes.
\end{itemize}

\subsection{Fluxos técnicos principais}
Descrever fluxos de sistema:
\begin{itemize}
    \item login e emissão de token;
    \item envio de redação;
    \item geração de feedback;
    \item consulta de redações;
    \item cálculo de analytics;
    \item carga de dados fictícios.
\end{itemize}

\section{Interpretação pedagógica}

\subsection{Competências do ENEM}
Descrever resumidamente C1, C2, C3, C4 e C5.

\begin{longtable}{|p{2cm}|p{5cm}|p{7cm}|}
\hline
\textbf{Competência} & \textbf{Nome curto} & \textbf{Descrição} \\
\hline
C1 & Norma escrita & preencher \\
\hline
C2 & Proposta e repertório & preencher \\
\hline
C3 & Argumentação & preencher \\
\hline
C4 & Coesão & preencher \\
\hline
C5 & Intervenção & preencher \\
\hline
\end{longtable}

\subsection{Cuidados na interpretação das notas}
Explicar que a avaliação automática deve apoiar, e não substituir, o julgamento pedagógico do professor.

\subsection{Uso em oficinas}
Explicar como usar os painéis para discussão com estudantes e professores.

\section{Administração e operação}

\subsection{Requisitos para execução}
Listar requisitos:
\begin{itemize}
    \item Python;
    \item Node.js;
    \item MongoDB;
    \item navegador;
    \item rede local, quando aplicável.
\end{itemize}

\subsection{Subida do backend}
Inserir comandos validados.

\subsection{Subida do frontend}
Inserir comandos para desenvolvimento e para oficina.

\subsection{Carga de dados fictícios}
Explicar o uso do script de seed.

\subsection{Health check da oficina}
Explicar como executar o script de verificação antes da oficina.

\subsection{Limpeza de dados por batch}
Explicar como limpar dados de uma oficina.

\section{Segurança e privacidade}

\subsection{Dados fictícios e dados reais}
Explicar a diferença e os cuidados necessários.

\subsection{Autenticação}
Descrever login, token e expiração em linguagem de manual.

\subsection{Perfis e permissões}
Explicar o que aluno e professor podem acessar.

\subsection{Recomendações para oficinas}
Listar cuidados:
\begin{itemize}
    \item usar dados fictícios;
    \item não expor MongoDB à internet;
    \item configurar CORS;
    \item usar chave JWT forte;
    \item limpar dados quando necessário.
\end{itemize}

\section{Solução de problemas}

\subsection{Backend não responde}
Descrever sintomas e verificações.

\subsection{Frontend mostra erro ou página 404}
Descrever rotas corretas e procedimentos.

\subsection{Login não funciona}
Descrever verificações de usuário, senha, backend e banco.

\subsection{Dados não aparecem}
Descrever verificações de carga, banco, batch e permissões.

\subsection{Analytics do professor não carrega}
Descrever verificações de turma, atividade, redações e cache.

\subsection{Erro de CORS}
Explicar como verificar `CORS_ORIGINS`.

\section{Perguntas frequentes}

\subsection{A nota automática substitui a correção do professor?}
Responder.

\subsection{Posso usar dados reais de estudantes?}
Responder com cuidados de consentimento, anonimização e segurança.

\subsection{Como reiniciar a aplicação?}
Responder.

\subsection{Como limpar os dados de uma oficina?}
Responder.

\subsection{Como criar novos usuários fictícios?}
Responder.

\section{Glossário}

\begin{longtable}{|p{4cm}|p{10cm}|}
\hline
\textbf{Termo} & \textbf{Definição} \\
\hline
Competência & preencher \\
\hline
Redação & preencher \\
\hline
Tema & preencher \\
\hline
Turma & preencher \\
\hline
Atividade & preencher \\
\hline
Seed batch & preencher \\
\hline
Analytics & preencher \\
\hline
\end{longtable}

\section{Apêndices}

\subsection{Apêndice A: usuários de demonstração}
Preencher com usuários gerados pelo script de carga.

\subsection{Apêndice B: comandos úteis}
Listar comandos de backend, frontend, MongoDB, carga e health check.

\subsection{Apêndice C: checklist pré-oficina}
Inserir checklist resumido.

\subsection{Apêndice D: histórico de versões do manual}

\begin{longtable}{|p{3cm}|p{4cm}|p{7cm}|}
\hline
\textbf{Data} & \textbf{Responsável} & \textbf{Alteração} \\
\hline
preencher & preencher & preencher \\
\hline
\end{longtable}

\end{document}
